\documentclass[11pt,a4paper]{article}
\usepackage[utf8]{inputenc}
\usepackage{amsmath,amssymb,amsthm}
\usepackage{algorithm}
\usepackage{algpseudocode}
\usepackage{graphicx}
\usepackage{hyperref}
\usepackage{booktabs}
\usepackage{bm}
\usepackage{bbm}

% Custom commands
\newcommand{\R}{\mathbb{R}}
\newcommand{\E}{\mathbb{E}}
\newcommand{\prob}{\mathbb{P}}
\newcommand{\vx}{\mathbf{x}}
\newcommand{\vy}{\mathbf{y}}
\newcommand{\vz}{\mathbf{z}}
\newcommand{\vtheta}{\boldsymbol{\theta}}
\newcommand{\sigmoid}{\sigma}
\DeclareMathOperator{\softmax}{softmax}

\title{Discrete Diffusion-by-Deblending EDA:\\
A Technical Documentation for the DbD-CS and DbD-CD Variants}
\author{Pateda Framework}
\date{\today}

\begin{document}

\maketitle

\begin{abstract}
This document provides a comprehensive mathematical description of the Discrete Diffusion-by-Deblending Estimation of Distribution Algorithm (DbD-EDA) for binary optimization problems. The DbD-EDA is a minimalist diffusion model that learns to transition between two distributions by training a neural network to predict the \emph{difference} (velocity) between blended samples and target samples. We focus on two main variants: DbD-CS (Current to Selected) and DbD-CD (Current to Closest in selected based on Distance). We describe the probabilistic blending process, the neural network architecture for deblending, the Gumbel-Softmax reparameterization, the learning procedure with multiple loss functions, the sampling mechanism, and the role of key hyperparameters including the number of alpha samples.
\end{abstract}

\tableofcontents
\newpage

%==============================================================================
\section{Introduction}
%==============================================================================

Estimation of Distribution Algorithms (EDAs) are population-based optimization methods that learn and sample from probabilistic models of promising solutions. The Discrete Diffusion-by-Deblending EDA (DbD-EDA) introduces a novel approach by leveraging concepts from flow-based diffusion models, adapted for binary optimization.

The key insight of DbD is that instead of learning to denoise corrupted data (as in traditional diffusion models), we learn the \emph{transformation} between two distributions: a source distribution $p_0$ and a target distribution $p_1$. The network learns to predict the ``velocity'' or difference that transforms samples from one distribution toward the other.

The key innovations of Discrete DbD-EDA include:
\begin{enumerate}
    \item \textbf{Probabilistic blending}: Creating intermediate samples by stochastically mixing two distributions using an interpolation parameter $\alpha$
    \item \textbf{Velocity prediction}: Training a neural network to predict the difference $(\vx_1 - \vx_0)$ given blended samples and the blending parameter $\alpha$
    \item \textbf{Gumbel-Softmax reparameterization}: Enabling differentiable discrete sampling during training
    \item \textbf{Iterative deblending}: Generating new samples by progressively applying the learned velocity field
    \item \textbf{Fitness-guided generation}: Optional conditioning on fitness values to bias sampling toward high-quality solutions
\end{enumerate}

%==============================================================================
\section{Problem Setting}
%==============================================================================

Consider a binary optimization problem where we seek to maximize (or minimize) an objective function $f: \{0,1\}^n \rightarrow \R$ over $n$ binary variables. In the context of an EDA, at each generation we have:

\begin{itemize}
    \item A \textbf{current population} $\mathcal{P}_{\text{current}} = \{\vx^{(1)}_c, \ldots, \vx^{(N)}_c\}$ representing the full population before selection
    \item A \textbf{selected population} $\mathcal{P}_{\text{selected}} = \{\vx^{(1)}_s, \ldots, \vx^{(M)}_s\}$ containing the top $M$ individuals by fitness, where $M = \lfloor \rho \cdot N \rfloor$ and $\rho$ is the selection ratio (truncation percent)
    \item Associated fitness values $\{f^{(1)}_c, \ldots, f^{(N)}_c\}$ and $\{f^{(1)}_s, \ldots, f^{(M)}_s\}$
\end{itemize}

The goal of DbD-EDA is to learn a generative model that can produce new candidate solutions that are likely to have high fitness by learning the transformation from one distribution to another.

%==============================================================================
\section{The Probabilistic Blending Process}
%==============================================================================

\subsection{Continuous vs. Discrete Blending}

In continuous domains, diffusion-by-deblending uses linear interpolation:
\begin{equation}
    \vx_\alpha = (1-\alpha)\vx_0 + \alpha \vx_1, \quad \alpha \in [0,1]
\end{equation}

For discrete (binary) variables, this deterministic interpolation is not directly applicable. Instead, we use \textbf{probabilistic blending}, where for each variable we sample from a mixture of two Bernoulli distributions.

\subsection{Probabilistic Blending for Binary Variables}

Given a pair of binary samples $\vx_0 \in \{0,1\}^n$ (from the source distribution) and $\vx_1 \in \{0,1\}^n$ (from the target distribution), the blending parameter $\alpha \in [0,1]$ defines the probability of selecting from each distribution.

For each variable $i$, the blended value is sampled as:
\begin{equation}
    p(x_{\alpha,i} = 1 | \vx_0, \vx_1, \alpha) = (1-\alpha) \cdot x_{0,i} + \alpha \cdot x_{1,i}
\end{equation}

This means:
\begin{equation}
    x_{\alpha,i} \sim \text{Bernoulli}\left((1-\alpha) \cdot x_{0,i} + \alpha \cdot x_{1,i}\right)
\end{equation}

\subsection{Interpretation}

The probabilistic blending creates a stochastic interpolation path:
\begin{itemize}
    \item At $\alpha = 0$: $x_{\alpha,i} = x_{0,i}$ (fully at source)
    \item At $\alpha = 1$: $x_{\alpha,i} = x_{1,i}$ (fully at target)
    \item At $\alpha = 0.5$: Equal probability of taking value from either distribution
\end{itemize}

\subsection{Creating Blended Training Samples}

The function \texttt{create\_blended\_binary\_samples} generates the training data:

\begin{algorithm}[H]
\caption{Create Blended Binary Samples for Training}
\label{alg:blending}
\begin{algorithmic}[1]
\Require Source population $P_0 = \{\vx^{(1)}_0, \ldots, \vx^{(N)}_0\}$
\Require Target population $P_1 = \{\vx^{(1)}_1, \ldots, \vx^{(N)}_1\}$
\Require Number of alpha samples $K$ per pair
\Require Fitness values $\{f^{(1)}_0, \ldots, f^{(N)}_0\}$ and $\{f^{(1)}_1, \ldots, f^{(N)}_1\}$ (optional)
\State Initialize empty lists: $\mathcal{A}$, $\mathcal{X}_{\alpha}$, $\mathcal{D}$, $\mathcal{F}$
\For{each pair $(\vx_0^{(j)}, \vx_1^{(j)})$ for $j = 1, \ldots, N$}
    \For{$k = 1$ to $K$}
        \State Sample $\alpha \sim \text{Uniform}(0, 1)$
        \State Compute mixture probability: $\mathbf{p} = (1-\alpha)\vx_0^{(j)} + \alpha \vx_1^{(j)}$
        \State Sample blended: $\vx_\alpha \sim \text{Bernoulli}(\mathbf{p})$ \Comment{Element-wise}
        \State Compute difference: $\mathbf{d} = \vx_1^{(j)} - \vx_0^{(j)}$ \Comment{Target for learning}
        \State Append $\alpha$ to $\mathcal{A}$
        \State Append $\vx_\alpha$ to $\mathcal{X}_{\alpha}$
        \State Append $\mathbf{d}$ to $\mathcal{D}$
        \If{fitness values provided}
            \State Compute blended fitness: $f_\alpha = (1-\alpha)f_0^{(j)} + \alpha f_1^{(j)}$
            \State Append $f_\alpha$ to $\mathcal{F}$
        \EndIf
    \EndFor
\EndFor
\State \Return $\mathcal{A}$, $\mathcal{X}_{\alpha}$, $\mathcal{D}$, $\mathcal{F}$
\end{algorithmic}
\end{algorithm}

%==============================================================================
\section{Neural Network Architecture}
%==============================================================================

\subsection{The Deblending Network}

The core component of DbD-EDA is a Multi-Layer Perceptron (MLP) called the \textbf{BinaryDeblendingNet}. Unlike traditional diffusion models that predict the clean data, this network learns to predict the \emph{difference} (velocity) between the target and source distributions.

\subsubsection{Network Structure}

The network architecture consists of:
\begin{enumerate}
    \item \textbf{Input layer}: Concatenation of:
        \begin{itemize}
            \item Blended binary data $\vx_\alpha \in \R^n$
            \item Blending parameter $\alpha \in \R$
            \item (Optional) Normalized fitness value $\tilde{f} \in \R$
        \end{itemize}
    \item \textbf{Hidden layers}: Sequence of fully-connected layers with configurable activation functions and dropout
    \item \textbf{Output layer}: Produces difference prediction $\hat{\mathbf{d}} \in \R^n$
\end{enumerate}

\subsubsection{Network Forward Pass}

Let $\vtheta$ denote the network parameters. The deblending network computes:

\begin{equation}
    \mathbf{h}_0 = \begin{cases}
        [\vx_\alpha; \alpha] \in \R^{n+1} & \text{if not fitness-guided} \\
        [\vx_\alpha; \alpha; \tilde{f}] \in \R^{n+2} & \text{if fitness-guided}
    \end{cases}
\end{equation}

\begin{equation}
    \mathbf{h}_\ell = \text{Dropout}\left(\phi_\ell\left(W_\ell \mathbf{h}_{\ell-1} + \mathbf{b}_\ell\right)\right), \quad \ell = 1, \ldots, L
\end{equation}

\begin{equation}
    \hat{\mathbf{d}} = W_{\text{out}} \mathbf{h}_L + \mathbf{b}_{\text{out}} \in \R^n
\end{equation}

where $\phi_\ell$ is the activation function for layer $\ell$ (configurable: ReLU, Tanh, Sigmoid, LeakyReLU, ELU, SELU, GELU, etc.).

\subsection{Network Inputs and Outputs}

\begin{table}[h]
\centering
\begin{tabular}{@{}lll@{}}
\toprule
\textbf{Component} & \textbf{Shape} & \textbf{Description} \\
\midrule
\textbf{Inputs:} & & \\
$\vx_\alpha$ & $(B, n)$ & Blended binary data (as floats) \\
$\alpha$ & $(B, 1)$ & Blending parameter \\
$\tilde{f}$ (optional) & $(B, 1)$ & Normalized fitness values \\
\midrule
\textbf{Outputs:} & & \\
$\hat{\mathbf{d}}$ & $(B, n)$ & Predicted difference $(\vx_1 - \vx_0)$ \\
\bottomrule
\end{tabular}
\caption{Network inputs and outputs, where $B$ is batch size and $n$ is number of variables.}
\label{tab:io}
\end{table}

\subsection{Interpretation of the Difference Prediction}

The network learns to predict:
\begin{equation}
    \hat{\mathbf{d}} = D_\vtheta(\vx_\alpha, \alpha) \approx \vx_1 - \vx_0
\end{equation}

For binary variables, the true difference $\mathbf{d} = \vx_1 - \vx_0$ takes values in $\{-1, 0, 1\}$:
\begin{itemize}
    \item $d_i = -1$: Bit flips from 1 to 0
    \item $d_i = 0$: Bit remains unchanged
    \item $d_i = 1$: Bit flips from 0 to 1
\end{itemize}

This formulation allows the network to learn the \emph{direction} of change needed to transform source samples toward the target distribution.

%==============================================================================
\section{Gumbel-Softmax Reparameterization}
%==============================================================================

\subsection{The Discrete Sampling Challenge}

During the sampling (deblending) process, we need to convert continuous probability values to discrete binary samples. However, standard discrete sampling (e.g., Bernoulli sampling) is non-differentiable, which poses challenges for gradient-based training.

\subsection{Gumbel-Softmax Distribution}

The Gumbel-Softmax trick provides a differentiable approximation to discrete sampling. Given logits $\mathbf{z} \in \R^{n \times 2}$ (representing probabilities for 0 and 1 for each variable) and temperature $\tau > 0$:

\begin{enumerate}
    \item Sample Gumbel noise:
    \begin{equation}
        g_{i,c} = -\log(-\log(u_{i,c})), \quad u_{i,c} \sim \text{Uniform}(0,1)
    \end{equation}

    \item Compute soft samples:
    \begin{equation}
        y_{i,c} = \frac{\exp((z_{i,c} + g_{i,c})/\tau)}{\sum_{c'=0}^{1} \exp((z_{i,c'} + g_{i,c'})/\tau)}
    \end{equation}
\end{enumerate}

\subsection{Straight-Through Estimator}

For obtaining hard (discrete) samples while maintaining gradient flow, the straight-through estimator is used:

\begin{equation}
    \hat{y}_{i,c} = \begin{cases}
        1 & \text{if } c = \arg\max_{c'} y_{i,c'} \\
        0 & \text{otherwise}
    \end{cases}
\end{equation}

\begin{equation}
    \tilde{y} = \hat{y} - \text{sg}(y) + y
\end{equation}

where $\text{sg}(\cdot)$ is the stop-gradient operator. This allows gradients to flow through $y$ in the backward pass while using $\hat{y}$ in the forward pass.

\subsection{Role in DbD-EDA Sampling}

In the sampling phase, the Gumbel-Softmax mechanism is used to:
\begin{enumerate}
    \item Convert the continuous probability predictions to binary samples
    \item Maintain stochasticity in the sampling process
    \item Allow for temperature-controlled exploration (higher temperature = more random, lower temperature = more deterministic)
\end{enumerate}

The sampling process uses Bernoulli sampling on the clipped probabilities:
\begin{equation}
    \vx_{\text{new}} = \text{Bernoulli}\left(\text{clip}(\vx + \Delta\alpha \cdot \hat{\mathbf{d}}, 0, 1)\right)
\end{equation}

%==============================================================================
\section{Learning Algorithm}
%==============================================================================

\subsection{Training Objective}

The model is trained to predict the difference between target and source samples given blended inputs. The primary objective minimizes the prediction error:

\begin{equation}
    \mathcal{L}(\vtheta) = \E_{\alpha, \vx_0, \vx_1, \vx_\alpha}\left[\ell\left(D_\vtheta(\vx_\alpha, \alpha), \vx_1 - \vx_0\right)\right]
\end{equation}

where $\ell$ is the loss function and the expectation is over:
\begin{itemize}
    \item Blending parameters $\alpha \sim \text{Uniform}(0, 1)$
    \item Source samples $\vx_0 \sim p_0$ (current or random population)
    \item Target samples $\vx_1 \sim p_1$ (selected population)
    \item Blended samples $\vx_\alpha \sim q(\vx_\alpha | \vx_0, \vx_1, \alpha)$
\end{itemize}

\subsection{Learning Procedure}

\begin{algorithm}[H]
\caption{Learning Binary DbD Model}
\label{alg:learning}
\begin{algorithmic}[1]
\Require Source population $P_0$, Target population $P_1$
\Require Fitness values $F_0$, $F_1$ (optional)
\Require Parameters: epochs $E$, batch size $B$, learning rate $\eta$, alpha samples $K$
\State Initialize deblending network $D_\vtheta$
\State Create blended samples: $(\mathcal{A}, \mathcal{X}_\alpha, \mathcal{D}, \mathcal{F}) \gets \texttt{create\_blended\_binary\_samples}(P_0, P_1, K, F_0, F_1)$
\For{epoch $= 1$ to $E$}
    \State Shuffle training data indices
    \For{each mini-batch $\mathcal{B}$ of size $B$}
        \State Extract batch: $\alpha_\mathcal{B}, \vx_{\alpha,\mathcal{B}}, \mathbf{d}_\mathcal{B}, f_\mathcal{B}$
        \State Forward pass: $\hat{\mathbf{d}}_\mathcal{B} \gets D_\vtheta(\vx_{\alpha,\mathcal{B}}, \alpha_\mathcal{B}, f_\mathcal{B})$
        \State Compute loss: $\mathcal{L} \gets \texttt{compute\_loss}(\hat{\mathbf{d}}_\mathcal{B}, \mathbf{d}_\mathcal{B}, f_\mathcal{B})$
        \State Backward pass: $\nabla_\vtheta \mathcal{L}$
        \State Update parameters: $\vtheta \gets \vtheta - \eta \nabla_\vtheta \mathcal{L}$
    \EndFor
\EndFor
\State \Return trained model $D_\vtheta$
\end{algorithmic}
\end{algorithm}

%==============================================================================
\section{Loss Functions}
%==============================================================================

The implementation supports multiple loss functions, each with different properties for guiding the learning process.

\subsection{Standard MSE Loss}

The default loss function measures the squared error between predicted and actual differences:

\begin{equation}
    \mathcal{L}_{\text{MSE}} = \frac{1}{Bn}\sum_{j=1}^{B}\sum_{i=1}^{n} \left(\hat{d}_i^{(j)} - d_i^{(j)}\right)^2
\end{equation}

where $\hat{d}_i^{(j)}$ is the predicted difference and $d_i^{(j)} = x_{1,i}^{(j)} - x_{0,i}^{(j)}$ is the true difference.

\subsection{Weighted MSE Loss}

Higher fitness samples receive higher weight in the loss:

\begin{equation}
    w^{(j)} = \gamma + (1 - \gamma) \cdot \frac{f^{(j)} - f_{\min}}{f_{\max} - f_{\min} + \epsilon}
\end{equation}

where $\gamma$ is the base weight (fitness\_weight parameter, default 0.1).

\begin{equation}
    \mathcal{L}_{\text{weighted}} = \frac{1}{Bn}\sum_{j=1}^{B} w^{(j)} \cdot \sum_{i=1}^{n}\left(\hat{d}_i^{(j)} - d_i^{(j)}\right)^2
\end{equation}

This encourages the model to more accurately predict transformations for high-fitness solutions.

\subsection{Ranking Loss}

The ranking loss prioritizes learning the relative ordering of solutions:

\begin{equation}
    \mathcal{L}_{\text{ranking}} = \mathcal{L}_{\text{MSE}} + \gamma \cdot \mathcal{L}_{\text{pairwise}}
\end{equation}

where the pairwise ranking component is:

\begin{equation}
    \mathcal{L}_{\text{pairwise}} = \frac{1}{|\mathcal{S}|} \sum_{(j,k) \in \mathcal{S}} L_\delta\left(\|\hat{\mathbf{d}}^{(j)}\|_2 - \|\hat{\mathbf{d}}^{(k)}\|_2 - \text{sign}(f^{(j)} - f^{(k)}) \cdot s\right)
\end{equation}

where:
\begin{itemize}
    \item $\mathcal{S}$ is a set of randomly sampled pairs
    \item $L_\delta$ is the Huber loss for robustness
    \item $s$ is a scaling factor (typically 0.1)
\end{itemize}

This loss encourages the network to preserve fitness ordering in its predictions.

\subsection{Huber Loss}

The Huber loss provides robustness to outliers:

\begin{equation}
    L_\delta(a) = \begin{cases}
        \frac{1}{2}a^2 & \text{if } |a| \leq \delta \\
        \delta\left(|a| - \frac{1}{2}\delta\right) & \text{otherwise}
    \end{cases}
\end{equation}

Applied element-wise:
\begin{equation}
    \mathcal{L}_{\text{Huber}} = \frac{1}{Bn}\sum_{j=1}^{B}\sum_{i=1}^{n} L_\delta\left(\hat{d}_i^{(j)} - d_i^{(j)}\right)
\end{equation}

%==============================================================================
\section{Fitness-Guided Generation}
%==============================================================================

\subsection{Conditional Deblending Network}

The fitness-guided variant conditions the deblending process on fitness information, inspired by conditional VAEs (C-VAE) and fitness-guided EDAs.

The network architecture is modified to include a fitness input:
\begin{equation}
    \tilde{f} = \frac{f - f_{\min}}{f_{\max} - f_{\min}}
\end{equation}

The concatenated input becomes:
\begin{equation}
    \mathbf{h}_0 = [\vx_\alpha; \alpha; \tilde{f}] \in \R^{n + 2}
\end{equation}

\subsection{Biasing Mechanism}

During training, the network learns the conditional distribution:
\begin{equation}
    D_\vtheta(\vx_\alpha, \alpha, \tilde{f}) \approx \E[\vx_1 - \vx_0 | \vx_\alpha, \alpha, f]
\end{equation}

This enables the network to:
\begin{enumerate}
    \item Learn that high-fitness solutions have certain structural patterns
    \item Associate different transformation directions with different fitness levels
    \item Generate solutions biased toward high-fitness regions when conditioned on high fitness values
\end{enumerate}

\subsection{Guided Sampling}

During sampling, we set a target fitness value (typically the maximum normalized fitness $\tilde{f} = 1.0$):
\begin{equation}
    \hat{\mathbf{d}} = D_\vtheta(\vx_\alpha, \alpha, \tilde{f}_{\text{target}} = 1.0)
\end{equation}

This biases the generation toward solutions that the model associates with high fitness.

%==============================================================================
\section{Sampling Algorithm}
%==============================================================================

\subsection{Iterative Deblending Process}

Sampling is performed via iterative application of the learned velocity field, progressively transforming samples from the initial distribution toward the target distribution.

\begin{algorithm}[H]
\caption{Sampling from Binary DbD Model}
\label{alg:sampling}
\begin{algorithmic}[1]
\Require Trained model $D_\vtheta$, number of samples $M$, denoising steps $S$, temperature $\tau$
\Require Initial population $P_{\text{init}}$ (for CS/CD variants) or random initialization
\State \textbf{Initialization:}
\If{DbD-CS variant}
    \State Sample $\vx_0$ from selected population with replacement
\ElsIf{DbD-CD variant}
    \State Sample $\vx_0$ from current population with replacement
\Else
    \State $\vx_0 \sim \text{Uniform}(\{0,1\}^n)$ or $\vx_0 \sim \text{Bernoulli}(0.5)^n$
\EndIf
\State Set $\vx \gets \vx_0$ as float tensor
\If{fitness-guided}
    \State Set target fitness: $\tilde{f}_{\text{target}} \gets 1.0$
\EndIf
\State Create alpha schedule: $\{\alpha_1, \alpha_2, \ldots, \alpha_S\}$ evenly spaced from $\frac{1}{S}$ to $1$
\State
\State \textbf{Iterative Deblending:}
\For{$s = 1$ to $S$}
    \State $\alpha_{\text{current}} \gets \alpha_{s-1}$ (with $\alpha_0 = 0$)
    \State $\alpha \gets \alpha_s$
    \State $\Delta\alpha \gets \alpha - \alpha_{\text{current}}$
    \State
    \If{fitness-guided}
        \State $\hat{\mathbf{d}} \gets D_\vtheta(\vx, \alpha, \tilde{f}_{\text{target}})$
    \Else
        \State $\hat{\mathbf{d}} \gets D_\vtheta(\vx, \alpha)$
    \EndIf
    \State
    \State \textbf{Update using velocity:}
    \State $\vx \gets \vx + \Delta\alpha \cdot \hat{\mathbf{d}}$
    \State $\vx \gets \text{clip}(\vx, 0, 1)$
    \State
    \State \textbf{Stochastic sampling:}
    \State $\mathbf{p} \gets \vx / \tau$ if $\tau \neq 1$ else $\vx$
    \State $\mathbf{p} \gets \text{clip}(\mathbf{p}, 0, 1)$
    \State $\vx \gets \text{Bernoulli}(\mathbf{p})$
\EndFor
\State \Return $\vx$ as integer binary samples
\end{algorithmic}
\end{algorithm}

\subsection{Progressive Transformation}

The key insight is the incremental update using the predicted velocity:
\begin{equation}
    \vx_{s+1} = \vx_s + \Delta\alpha_s \cdot D_\vtheta(\vx_s, \alpha_s)
\end{equation}

As $s$ increases (and thus $\alpha$ increases from 0 to 1):
\begin{itemize}
    \item The samples progressively move toward the target distribution
    \item Early steps make larger structural changes
    \item Later steps refine the solutions
    \item The Bernoulli sampling at each step maintains discrete structure
\end{itemize}

%==============================================================================
\section{DbD-CS vs DbD-CD Variants}
%==============================================================================

The two main variants differ in how the source population $P_0$ is constructed and how the target samples are paired.

\subsection{DbD-CS: Current to Selected}

In the DbD-CS (Current to Selected) variant:

\begin{itemize}
    \item \textbf{Source population ($P_0$)}: The current population before selection (or a random population in the first generation)
    \item \textbf{Target population ($P_1$)}: The selected population (top individuals by fitness)
    \item \textbf{Pairing}: Random pairing between source and target samples
    \item \textbf{Sampling initialization}: Samples are initialized from the selected population
\end{itemize}

\textbf{Learning:}
\begin{equation}
    D_\vtheta \text{ learns: } \vx_{\text{current}} \rightarrow \vx_{\text{selected}}
\end{equation}

\textbf{Characteristics:}
\begin{itemize}
    \item Learns the general direction of fitness improvement
    \item May learn longer-range transformations
    \item Better for exploration when current and selected populations are diverse
\end{itemize}

\subsection{DbD-CD: Current to Closest in selected (Distance-based)}

In the DbD-CD (Current to Closest by Distance) variant:

\begin{itemize}
    \item \textbf{Source population ($P_0$)}: The current population before selection
    \item \textbf{Target population ($P_1$)}: For each source sample, the \emph{closest} sample in the selected population based on Hamming distance
    \item \textbf{Pairing}: Each source sample is paired with its nearest neighbor in the selected population
    \item \textbf{Sampling initialization}: Samples are initialized from the current population
\end{itemize}

\textbf{Closest Neighbor Computation:}
\begin{equation}
    \vx_1^{(j)} = \arg\min_{\vx \in P_{\text{selected}}} d_H(\vx_0^{(j)}, \vx)
\end{equation}
where $d_H(\cdot, \cdot)$ is the Hamming distance.

\textbf{Learning:}
\begin{equation}
    D_\vtheta \text{ learns: } \vx_{\text{current}} \rightarrow \text{closest}(\vx_{\text{current}}, P_{\text{selected}})
\end{equation}

\textbf{Characteristics:}
\begin{itemize}
    \item Learns smaller, more local transformations
    \item Focuses on minimal changes needed for improvement
    \item Better for exploitation and fine-tuning
    \item Reduces the complexity of the learning task
\end{itemize}

\subsection{Comparison Summary}

\begin{table}[h]
\centering
\begin{tabular}{@{}lcc@{}}
\toprule
\textbf{Aspect} & \textbf{DbD-CS} & \textbf{DbD-CD} \\
\midrule
Source ($P_0$) & Current/Random population & Current population \\
Target ($P_1$) & Selected population & Closest in selected \\
Pairing & Random & Nearest neighbor \\
Sampling init & From selected & From current \\
Transformation & Long-range & Local \\
Emphasis & Exploration & Exploitation \\
\bottomrule
\end{tabular}
\caption{Comparison of DbD-CS and DbD-CD variants.}
\label{tab:variants}
\end{table}

%==============================================================================
\section{Role of \texttt{num\_alpha\_samples}}
%==============================================================================

The parameter \texttt{num\_alpha\_samples} (denoted $K$) controls how many blended training samples are created for each pair of source and target samples.

\subsection{Training Data Augmentation}

For each pair $(\vx_0^{(j)}, \vx_1^{(j)})$, the algorithm creates $K$ training examples:
\begin{equation}
    \{(\alpha_k, \vx_{\alpha_k}^{(j)}, \mathbf{d}^{(j)})\}_{k=1}^{K}
\end{equation}

where each $\alpha_k \sim \text{Uniform}(0, 1)$ is independently sampled.

\subsection{Impact on Learning}

\textbf{Total training samples:}
\begin{equation}
    N_{\text{train}} = N \times K
\end{equation}

where $N$ is the number of source-target pairs.

\textbf{Effects of increasing $K$:}
\begin{enumerate}
    \item \textbf{Better coverage of $\alpha$ values}: More alpha samples means the network sees the transformation at more intermediate stages
    \item \textbf{Improved velocity field estimation}: Denser sampling of the interpolation path leads to more accurate velocity predictions
    \item \textbf{Reduced overfitting}: More training data with variety helps generalization
    \item \textbf{Computational cost}: Linear increase in training time and memory
\end{enumerate}

\subsection{Recommended Values}

\begin{itemize}
    \item \textbf{Default}: $K = 20$ (a good balance between quality and efficiency)
    \item \textbf{Small populations}: Consider $K = 30$--$50$ to compensate for fewer pairs
    \item \textbf{Large populations}: $K = 10$--$20$ may suffice
    \item \textbf{Complex problems}: Higher $K$ can help learn intricate transformations
\end{itemize}

\subsection{Mathematical Interpretation}

The parameter $K$ controls the granularity of the learned velocity field. With more alpha samples, the network learns a more continuous approximation to:

\begin{equation}
    D_\vtheta(\vx_\alpha, \alpha) \approx \frac{\partial \vx_\alpha}{\partial \alpha} = \vx_1 - \vx_0
\end{equation}

This is analogous to numerical differentiation: more samples lead to better approximation of the derivative (velocity) along the interpolation path.

%==============================================================================
\section{Complete EDA Algorithm}
%==============================================================================

\begin{algorithm}[H]
\caption{Discrete DbD-EDA (CS/CD Variants)}
\label{alg:eda}
\begin{algorithmic}[1]
\Require Fitness function $f$, population size $N$, generations $G$, selection ratio $\rho$
\Require DbD parameters: hidden\_dims, epochs, batch\_size, num\_alpha\_samples $K$, n\_steps $S$
\Require Variant: \texttt{dbd\_cs} or \texttt{dbd\_cd}
\State
\State \textbf{Initialization:}
\State $\mathcal{P}_0 \gets$ random binary population of size $N$
\State Evaluate: $f^{(i)} \gets f(\vx^{(i)})$ for all $\vx^{(i)} \in \mathcal{P}_0$
\State Track best: $\vx^* \gets \arg\max_{\vx \in \mathcal{P}_0} f(\vx)$
\State $P_{\text{prev}} \gets \texttt{None}$
\State
\For{generation $g = 1$ to $G$}
    \State \textbf{Selection:}
    \State $M \gets \lfloor \rho \cdot N \rfloor$
    \State $\mathcal{S}_g \gets$ top-$M$ individuals from $\mathcal{P}_{g-1}$ by fitness
    \State
    \State \textbf{Prepare Source Population:}
    \If{$P_{\text{prev}}$ is \texttt{None}}
        \State $P_0 \gets$ random binary samples of size $M$
    \Else
        \State $P_0 \gets$ sample $M$ from $P_{\text{prev}}$
    \EndIf
    \State
    \State \textbf{Prepare Target Population:}
    \If{variant is \texttt{dbd\_cs}}
        \State $P_1 \gets \mathcal{S}_g$
    \ElsIf{variant is \texttt{dbd\_cd}}
        \State $P_1 \gets \texttt{find\_closest\_neighbors}(P_0, \mathcal{S}_g)$
    \EndIf
    \State
    \State \textbf{Learning:}
    \State $D_\vtheta \gets \texttt{learn\_binary\_dbd}(P_0, P_1, \text{params})$
    \State
    \State \textbf{Sampling:}
    \If{variant is \texttt{dbd\_cs}}
        \State $\mathcal{P}_g \gets \texttt{sample\_binary\_dbd\_cs}(D_\vtheta, N, \mathcal{S}_g, \text{sampling\_params})$
    \ElsIf{variant is \texttt{dbd\_cd}}
        \State $\mathcal{P}_g \gets \texttt{sample\_binary\_dbd\_cd}(D_\vtheta, N, \mathcal{P}_{g-1}, \text{sampling\_params})$
    \EndIf
    \State
    \State \textbf{Evaluation:}
    \State $f^{(i)} \gets f(\vx^{(i)})$ for all $\vx^{(i)} \in \mathcal{P}_g$
    \State
    \State \textbf{Update best:}
    \If{$\max_{\vx \in \mathcal{P}_g} f(\vx) > f(\vx^*)$}
        \State $\vx^* \gets \arg\max_{\vx \in \mathcal{P}_g} f(\vx)$
    \EndIf
    \State
    \State $P_{\text{prev}} \gets \mathcal{P}_{g-1}$
    \State
    \State \textbf{Early termination:}
    \If{$|f(\vx^*) - f_{\text{optimal}}| < \epsilon$}
        \State \textbf{break}
    \EndIf
\EndFor
\State \Return $\vx^*$, $f(\vx^*)$
\end{algorithmic}
\end{algorithm}

%==============================================================================
\section{Summary of Mathematical Formulations}
%==============================================================================

\subsection{Probabilistic Blending}
\begin{equation}
    x_{\alpha,i} \sim \text{Bernoulli}\left((1-\alpha) x_{0,i} + \alpha x_{1,i}\right), \quad \forall i \in \{1, \ldots, n\}
\end{equation}

\subsection{Deblending Network}
\begin{equation}
    \hat{\mathbf{d}} = D_\vtheta(\vx_\alpha, \alpha) \approx \vx_1 - \vx_0
\end{equation}

\subsection{Loss Functions}
\begin{align}
    \mathcal{L}_{\text{MSE}} &= \frac{1}{Bn}\sum_{j,i} \left(\hat{d}_i^{(j)} - d_i^{(j)}\right)^2 \\
    \mathcal{L}_{\text{weighted}} &= \frac{1}{Bn}\sum_{j,i} w^{(j)} \left(\hat{d}_i^{(j)} - d_i^{(j)}\right)^2 \\
    \mathcal{L}_{\text{ranking}} &= \mathcal{L}_{\text{MSE}} + \gamma \cdot \mathcal{L}_{\text{pairwise}} \\
    \mathcal{L}_{\text{Huber}} &= \frac{1}{Bn}\sum_{j,i} L_\delta\left(\hat{d}_i^{(j)} - d_i^{(j)}\right)
\end{align}

\subsection{Sampling Update}
\begin{equation}
    \vx_{s+1} = \text{Bernoulli}\left(\text{clip}\left(\vx_s + \Delta\alpha_s \cdot D_\vtheta(\vx_s, \alpha_s), 0, 1\right)\right)
\end{equation}

\subsection{Fitness Normalization}
\begin{equation}
    \tilde{f} = \frac{f - f_{\min}}{f_{\max} - f_{\min} + \epsilon}
\end{equation}

\subsection{Hamming Distance (for DbD-CD)}
\begin{equation}
    d_H(\vx, \vy) = \sum_{i=1}^{n} \mathbbm{1}[x_i \neq y_i]
\end{equation}

%==============================================================================
\section{Conclusion}
%==============================================================================

The Discrete Diffusion-by-Deblending EDA provides a powerful framework for binary optimization by:

\begin{enumerate}
    \item Leveraging probabilistic blending to create smooth interpolations between distributions
    \item Training a neural network to learn the velocity field for transforming source to target distributions
    \item Using Gumbel-Softmax for differentiable discrete sampling
    \item Supporting multiple loss functions for different optimization scenarios
    \item Enabling fitness-guided generation to bias toward high-quality solutions
    \item Providing variant-specific behaviors (exploration via CS, exploitation via CD)
\end{enumerate}

The \texttt{num\_alpha\_samples} parameter ($K$) plays a crucial role in determining the quality of the learned velocity field by controlling the density of training examples along the interpolation path. Higher values of $K$ lead to better approximation of the transformation but at increased computational cost.

The algorithm is particularly suited for problems where:
\begin{itemize}
    \item Variable interactions are complex and difficult to model explicitly
    \item The fitness landscape has multiple local optima
    \item A balance between exploration and exploitation is critical
    \item The transformation between current and selected populations encodes useful fitness-improving patterns
\end{itemize}

\end{document}
